\section[Introduction]{Introduction and Chapter Scope}
\label{sec:pd_intro}

The relentless growth of data-centre interconnect capacity has driven the O-band
silicon-photonics ecosystem toward 400G and beyond. IEEE~802.3 Clause~140
specifies 400GBASE-DR4 as a four-lane PAM-4 interface at \SI{1310}{\nano\metre},
placing each lane at \SI{106.25}{\giga\bit\per\second} raw line rate and demanding
a receiver-side photodetector that is simultaneously fast, efficient, and quiet.
Meeting that specification on a CMOS-foundry platform is not merely a component
selection problem: it requires a design methodology that connects photon absorption
physics, material constraints, device topology, and electrode engineering into a
single coherent argument.

This chapter develops that argument from first principles. The discussion is
structured as a progressive narrowing: beginning with the electromagnetic
foundations of light absorption in semiconductors, passing through the metrics
that define a useful receiver device, surveying the landscape of competing
architectures and materials, and arriving at a specific vertical n-i-p
Ge-on-Si waveguide photodetector whose geometry, contact topology, and
simulation strategy are each justified by the physics established along the way.

The chapter is organised as follows.
\Cref{sec:pd_fundamentals} establishes the physics of photodetection, deriving
the Beer-Lambert absorption law from Maxwell's equations, developing the
drift-diffusion carrier transport framework, and quantifying the recombination
mechanisms that limit carrier collection.
\Cref{sec:pd_metrics} defines the evaluation language, namely responsivity,
bandwidth, dark current, noise, and the trade-offs among them, that governs
every architectural and material choice that follows.
\Cref{sec:pd_architectures} surveys photodetector device classes from the
photoconductor through the avalanche photodiode and derives, rather than
asserts, the selection of the waveguide-integrated PIN junction for high-speed
integrated links.
\Cref{sec:pd_material} justifies germanium-on-silicon as the material platform,
including its spectral suitability at \SI{1310}{\nano\metre}, its CMOS integration
path, and the threading-dislocation challenge that sets the dark-current floor.
\Cref{sec:pd_literature} reviews the experimental record of Ge-on-Si waveguide
photodetectors critically and identifies the gap that motivates the U-shaped
electrode approach adopted here.
\Cref{sec:pd_design} presents the proposed device architecture in detail,
including the vertical PIN layer stack, the adiabatic taper coupling strategy,
the U-shaped electrode topology, and the geometry-to-performance design space
mapping.
\Cref{sec:pd_simulation} describes the multiphysics simulation strategy: FDTD
optical modelling, CHARGE drift-diffusion solving, MATLAB post-processing, and
the validation approach anchoring the simulation to the published benchmark.
\Cref{sec:pd_summary} closes with a synthesis of the key design insights and a
transition to the results chapter.
